\section{Certificate scope, observers, and canonical export}
\label{sec:certificates}

Certificates are unavoidable in constructive mathematics: factorizations,
paths, finite covers, branch choices, approximation traces, and bounded-search
witnesses all carry information.  The mistake is to make every internal
certificate part of the public theorem.  DARC performs certificate escape
analysis within the same demand lattice used for precision and residuals.

\subsection{Scoped certificates}

A scoped certificate contract over an outer context \(\Gamma\) and public goal
\(\Phi\) consists of:
\begin{enumerate}
  \item a validity formula \(V(c)\) in a scope with one additional variable;
  \item a consumer derivation
        \(V(c),\Gamma^+\vdash\Phi^+\).
\end{enumerate}
A producer proves \(\Gamma\vdash\exists c\,V(c)\).  By
\cref{lem:exseal}, the certificate is sealed.

\begin{theorem}[Local certificate erasure]
\label{thm:local-cert-erasure}
If no demanded public formula or live observer contains the certificate
variable, then a valid local certificate can be eliminated without selecting a
canonical code.
\end{theorem}

\begin{proof}
The demand condition ensures that all active uses of the certificate occur
inside its consumer subgraph and that the subgraph's exported conclusion is in
the outer scope.  Compile that subgraph to
\(V(c),\Gamma^+\vdash\Phi^+\), then apply existential elimination to the
producer.  Since the conclusion is scoped outside the binder, the certificate
cannot occur in it.
\end{proof}

This is stronger and simpler than finding a least certificate.  In particular,
nonuniqueness causes no difficulty when only an invariant consequence is
needed.

\subsection{Observed certificates}

Sometimes the public result depends on an observation \(O(c)\) but not on the
full certificate.  There are three common compatibility modes.

\paragraph{Exact stability.}
For all valid \(c,c'\), prove \(O(c)=O(c')\).  One may seal a local
certificate, define the public value by any valid witness inside the proof,
and use stability to show independence.

\paragraph{Quotient stability.}
Prove an explicit equivalence relation \(E(O(c),O(c'))\).  The public theorem
states a quotient-invariant property or uses a canonical finite code for the
equivalence class only when such a code has a bounded constructive
normalization.

\paragraph{Metric stability.}
At requested precision \(p\), prove
\(\operatorname{dist}(O(c),O(c'))\le 2^{-p}\) or another explicit tolerance.
The observation is then well-defined at that finite precision.  A global exact
object requires additional coherence and a tail modulus; DARC does not infer
one automatically.

\begin{theorem}[Stable observer elimination]
\label{thm:observer-elimination}
Let \(V(c)\) be a local certificate predicate and let \(P(z)\) be a public
property.  Suppose:
\begin{enumerate}
  \item \(\Gamma\vdash\exists c\,V(c)\);
  \item \(V(c),\Gamma^+\vdash P(O(c))\);
  \item the public conclusion is invariant under the chosen exact, quotient,
        or finite-precision compatibility relation.
\end{enumerate}
Then the certificate may remain scoped, and only the observation contract is
exported.
\end{theorem}

\begin{proof}
Open the existential certificate.  Derive the observed property from the local
consumer.  Use the compatibility theorem whenever the public statement
compares this observation with one obtained through another local witness.
Finally close the existential.  No representative is selected outside the
scope.
\end{proof}

The theorem is a schema because the exact public formula differs by
compatibility mode.  Its logical core is always existential elimination plus
an ordinary compatibility derivation.

\subsection{Named certificate export by bounded least search}

A public theorem may genuinely require a witness: for example, a least prime
factor, a smallest repeated index, a canonical finite path, or a literal code
used by a downstream endpoint.  Then canonicalization is appropriate, but it
requires two pieces of constructive data:
\begin{enumerate}
  \item a finite bound \(N(x)\);
  \item a fixed primitive-recursive validity bit \(f(i,x)\), with
        \(f(i,x)=0\) reflecting validity.
\end{enumerate}
The checked search symbol \(\mu_f(N,x)\) returns the least zero below \(N\), or
\(N\) if none exists.

\begin{theorem}[Named export]
\label{thm:named-export}
Assume
\[
  \Gamma\vdash f(i,x)=0,
  \qquad
  \Gamma\vdash i<N.
\]
Then \(\Gamma\) proves
\begin{align*}
  \mu_f(N,x)&<N,\\
  f(\mu_f(N,x),x)&=0,\\
  \mu_f(N,x)&\le i.
\end{align*}
Consequently the least valid code is a legitimate public term.
\end{theorem}

\begin{proof}
This is the checked theorem \code{Search.search\_found}.  Its proof uses the
primitive-recursive search recurrence and one quantifier-free induction.  The
DARC wrapper merely exposes it as the named-certificate operation.
\end{proof}

The theorem requires existence below a bound.  It does not turn unbounded
existence into a witness.  A height bound or finite enumeration theorem must be
proved or left residual.

\subsection{Escape analysis as a demand transfer}

For each certificate field \(c_j\), introduce a Boolean coordinate
\(\Cert(c_j)\).  A public field occurrence seeds that coordinate.  An observer
transfer may replace demand for the full certificate by demand for a stable
observation.  A named-export transfer adds the validity checker, bound, and any
guards needed by canonicalization.  If no rule reaches \(\Cert(c_j)\), the
field is dead and must not occur in the compiled endpoint.

This analysis is field-sensitive.  A certificate may contain:
\[
  c=(\text{trace},\text{branch label},\text{error bound},\text{provenance}).
\]
A value query may need only the error bound; a replay query may need the trace;
a labeled-root query may need the branch and separation guard.  Treating the
certificate as one indivisible code would retain unnecessary information and
residuals.

\subsection{Certificate fusion}

Suppose two adjacent nodes produce certificates \(c\) and \(d\), and the
consumer uses only a relation \(R(c,d)\) and a stable observation.  Rather than
naming both separately, form one scoped certificate
\(\exists c,d\,[V(c)\wedge W(d)\wedge R(c,d)]\) and eliminate it locally.
This is particularly useful for finite relational runs.  It prevents the
invalid move of choosing a valid \(c\) and a valid \(d\) independently when no
compatible pair exists.

\subsection{No hidden choice}

The three export modes make choice obligations visible:
\begin{itemize}
  \item sealed mode uses no choice beyond an existential proof;
  \item observed mode uses a compatibility theorem;
  \item named mode uses bounded least search.
\end{itemize}
There is no fourth mode that silently extracts an arbitrary witness from an
unbounded existential.  This restriction is not bureaucratic; it is what keeps
the compiled proof intuitionistic and primitive recursive.
